\chapter[Introdução]{Introdução}

A engenharia aplicada à saúde tem desempenhado um papel fundamental no desenvolvimento de soluções inovadoras para o diagnóstico e monitoramento 
de doenças cardiovasculares. Dentre essas doenças, as arritmias cardíacas se destacam pela alta prevalência e pelo potencial risco à vida, 
tornando sua detecção e classificação temas de grande importância clínica e científica. O eletrocardiograma (ECG) permanece como principal 
exame para análise da atividade elétrica cardíaca, mas desafios relacionados ao grande volume de dados, à subjetividade e à possibilidade de 
erros na interpretação manual evidenciam a necessidade de métodos automatizados e precisos.

Neste contexto, abordagens baseadas em inteligência artificial, especialmente redes neurais convolucionais, têm sido exploradas com sucesso para 
a classificação automática de batimentos cardíacos, contribuindo para maior agilidade, acurácia e confiabilidade nos diagnósticos. A presente 
pesquisa propõe o desenvolvimento e a comparação de modelos CNN1D e CNN2D aplicados à identificação e classificação de arritmias em sinais de 
ECG, tendo como base de dados o repositório MIT-BIH, referência mundial na área.

Nos subtópicos seguintes serão apresentados a contextualização do tema, a justificativa para a realização deste trabalho, os objetivos 
estabelecidos e a estrutura adotada para a elaboração deste trabalho.

\section{Contextualização}

As arritmias cardíacas constituem um importante problema de saúde pública, associando-se a maior risco de hospitalizações, morbidade e mortalidade 
cardiovascular. Estudos recentes sobre a epidemiologia dos transtornos de condução e das arritmias cardíacas no Brasil apontam mais de 700 mil 
internações entre 2014 e 2024, com aumento progressivo dos casos e predominância de atendimentos em caráter de urgência, o que sugere diagnóstico 
tardio e manejo muitas vezes em situações críticas \citeonline{Brito}. Análises complementares mostram que esses eventos acometem com maior 
frequência indivíduos idosos, especialmente acima de 80 anos, sobrecarregando o sistema de saúde e reforçando a necessidade de estratégias de 
detecção precoce \citeonline{Jesus}.

A Sociedade Brasileira de Cardiologia destaca o eletrocardiograma de superfície como exame fundamental para avaliação do ritmo cardíaco, 
por ser um método não invasivo, amplamente disponível e de baixo custo, essencial para o diagnóstico de arritmias em diversos níveis de 
atenção à saúde \citeonline{Samesima}. Porém, a interpretação manual dos registros de ECG pode ser dificultada pela variabilidade dos batimentos, 
pela presença de artefatos e pelo volume de dados gerados em contextos de monitorização contínua, aumentando 
o risco de erros.

Nesse cenário, abordagens baseadas em inteligência artificial têm ganhado destaque como ferramentas de apoio ao diagnóstico. Revisões de 
literatura sobre modelos de deep learning para detecção de arritmias em ECG mostram que redes neurais convolucionais aplicadas à base MIT‑BIH 
Arrhythmia Database atingem acurácia superior a 96\% na classificação de batimentos cardíacos, reforçando o potencial dessas técnicas 
para aprimorar a detecção automática de arritmias \citeonline{Ansari} e \citeonline{Cretu}. Dessa forma, a integração entre sinais de ECG e 
modelos computacionais avançados surge como alternativa promissora para sistemas de monitoramento cardíaco.

\section{Justificativa}

As arritmias cardíacas geram grande impacto clínico e econômico, com elevado número de internações, uso de recursos especializados e risco 
aumentado de eventos fatais, o que torna necessário aprimorar estratégias de detecção e acompanhamento desses pacientes. A carga de hospitalizações 
por transtornos de condução e arritmias no Brasil evidencia uma demanda crescente por soluções que auxiliem o diagnóstico de forma mais rápida, 
padronizada e acessível, especialmente em serviços com alta demanda e oferta limitada de especialistas. O ECG é de fácil acesso e apresenta bons
resultados, porém sempre depende de um especialista para dar o diagnóstico, sendo que às vezes esse diagnóstico pode ser equivocado e demorado,
então surge como alternativa a utilização de modelos para a detecção e classificação de arritmias cardíacas. 

\section{Objetivos}

O objetivo geral do trabalho é desenvolver dois modelos de redes neurais convolucionais (CNN), sendo um modelo de CNN1D e um modelo CNN2D, 
utilizando a base de dados pública e gratuita do MIT-BIH, visando detectar e classificar os batimentos com arritmia. Para alcançar o objetivo geral 
é necessário coletar os dados, fazer o pré-processamento, separar de acordo com a classificação dos especialistas, treinar os modelos, verificar
as métricas e comparar com resultados já obtidos por outros trabalhos. 

\section{Estrutura do Trabalho}

Este trabalho está organizado em seis capítulos, dispostos da seguinte forma:

O \textbf{Capítulo 2} apresenta a \textbf{Fundamentação Teórica}, abordando os conceitos essenciais para a compreensão do estudo, incluindo o 
funcionamento das CNNs, técnicas de regularização, otimizadores e as métricas de avaliação de desempenho utilizadas para a classificação de arritmias.

O \textbf{Capítulo 3} descreve os \textbf{Trabalhos Relacionados}, contextualizando a pesquisa no estado da arte atual. Nesta seção, são discutidas 
abordagens anteriores que utilizaram modelos unidimensionais e bidimensionais, estabelecendo um comparativo direto com a metodologia proposta.

O \textbf{Capítulo 4} detalha a \textbf{Metodologia} adotada, descrevendo a base de dados MIT-BIH Arrhythmia Database, as etapas de pré-processamento 
dos sinais (incluindo normalização e segmentação), a geração de espectrogramas via STFT e a arquitetura das redes CNN1D e CNN2D desenvolvidas.

O \textbf{Capítulo 5} expõe os \textbf{Resultados e Discussão}, apresentando as métricas de desempenho obtidas (acurácia, sensibilidade, especificidade, 
precisão e F1-Score), as matrizes de confusão e uma análise comparativa entre as abordagens temporal (1D) e espectro-temporal (2D).

Por fim, o \textbf{Capítulo 6} apresenta a \textbf{Conclusão}, sintetizando as contribuições do trabalho, as limitações encontradas e sugerindo 
direções para trabalhos futuros.